\section{幾何による別解}\label{節:幾何による別解}
微分を学ぶと,最大値や接線を同じ手順で求められるようになる.ただし,図形の性質から直接求まることもある.二つの方法を比べてみよう.

\begin{example}{円に内接する長方形}
半径$1$の円に内接する長方形の面積の最大値を求めよ.
\end{example}
円の中心を原点とし,長方形の辺が座標軸に平行になるように取る.第1象限の頂点を$(x,y)$とすれば,\,$x,y>0$,\,$x^2+y^2=1$であり,面積は$S=4xy$である.\,$y=\sqrt{1-x^2}$と表して微分すると,
\[
 S(x)=4x\sqrt{1-x^2},\qquad
 S'(x)=\frac{4(1-2x^2)}{\sqrt{1-x^2}}.
\]
$0<x<1$では分母は正なので,\,$x=\frac1{\sqrt2}$の前後で増加から減少へ変わる.このとき$y=\frac1{\sqrt2}$であり,最大値は$2$となる.

\noindent\textbf{【別解】}\quad
同じ式を$x,y$のまま見てみよう.平方が非負であることから,
\[
 (x-y)^2\ge0,\qquad 2xy\le x^2+y^2=1.
\]
したがって$S=4xy\le2$である.等号は$x=y=\frac1{\sqrt2}$のときに成り立つ.図形では,長方形が正方形になる場合である.

幾何的にも,長方形の辺を$a,b$,対角線を$2$とすれば,$a,b$を二辺とする直角三角形の斜辺は$2$である.この直角三角形を四つ並べ,外側を一辺$2$の正方形,中央を一辺$|a-b|$の正方形にすると,面積の和から
\[
 4=4\cdot\frac{ab}{2}+(a-b)^2
\]
となる.したがって$ab\le2$であり,等号は中央の正方形の面積が$0$,すなわち$a=b$のときに成り立つ.

微分による解法は,円を別の曲線へ変えても,一つの変数で面積を表せれば使える.一方,平方を使った解法では,円の条件$x^2+y^2=1$がそのまま上からの評価になる.答えの正方形も,等号条件$x=y$から分かる.このように,計算の長さだけでなく,条件のどの部分を使ったかまで比べておきたい.

